% Môn: Vật lý
% Lớp: 12
% Chương: Chương I. Vật lí nhiệt
% Bài: L12C1 Bài 1. Cấu trúc của chất. Sự chuyển thể · Giải thích hiện tượng và ứng dụng liên quan đến sự hóa hơi
% ===== Câu 41 =====
% ID: L12C1B1-52-DS
% Mức: NB
\dangbt{Giải thích hiện tượng và ứng dụng liên quan đến sự hóa hơi}
\begin{ex}
% ID: L12C1B1-52-DS
% Mức: NB
Xét máy lạnh.
\choiceTF
{\True Môi chất bay hơi ở dàn lạnh và thu nhiệt.}
{\True Môi chất ngưng tụ ở dàn nóng và tỏa nhiệt.}
{\True Máy lạnh chuyển nhiệt từ nơi lạnh sang nơi nóng nhờ công của máy nén.}
{Máy lạnh tạo ra năng lượng từ hư không.}
\loigiai{
\begin{itemchoice}
\item (Đúng) Môi chất bay hơi ở dàn lạnh và thu nhiệt.
\item (Đúng) Môi chất ngưng tụ ở dàn nóng và tỏa nhiệt.
\item (Đúng) Máy lạnh chuyển nhiệt từ nơi lạnh sang nơi nóng nhờ công của máy nén.
\item (Sai) Máy lạnh tạo ra năng lượng từ hư không.
\end{itemchoice}
}
\end{ex}

% ===== Câu 45 =====
% ID: L12C1B1-56-DS
% Mức: TH
\begin{ex}
% ID: L12C1B1-56-DS
% Mức: TH
Xét nồi áp suất.
\choiceTF
{\True Áp suất trong nồi lớn hơn khí quyển.}
{\True Nhiệt độ sôi của nước tăng.}
{\True Thức ăn có thể chín nhanh hơn do nhiệt độ nấu cao.}
{Mở nắp ngay khi còn áp suất cao là an toàn.}
\loigiai{
\begin{itemchoice}
\item (Đúng) Áp suất trong nồi lớn hơn khí quyển.
\item (Đúng) Nhiệt độ sôi của nước tăng.
\item (Đúng) Thức ăn có thể chín nhanh hơn do nhiệt độ nấu cao.
\item (Sai) Mở nắp ngay khi còn áp suất cao là an toàn.
\end{itemchoice}
}
\end{ex}

% ===== Câu 49 =====
% ID: L12C1B1-67-DS
% Mức: VD
\begin{ex}
% ID: L12C1B1-67-DS
% Mức: VD
Xét hiện tượng bay hơi trong đời sống.
\choiceTF
{\True Bay hơi xảy ra ở mặt thoáng.}
{\True Nhiệt độ cao thường làm bay hơi nhanh hơn.}
{\True Gió có thể làm quần áo khô nhanh hơn.}
{Bay hơi luôn chỉ xảy ra ở 100 °C.}
\loigiai{
\begin{itemchoice}
\item (Đúng) Bay hơi xảy ra ở mặt thoáng.
\item (Đúng) Nhiệt độ cao thường làm bay hơi nhanh hơn. (Vì): ay hơi có thể xảy ra ở mọi nhiệt độ; ba yếu tố thường gặp là nhiệt độ, diện tích mặt thoáng và gió
\item (Đúng) Gió có thể làm quần áo khô nhanh hơn.
\item (Sai) Bay hơi luôn chỉ xảy ra ở 100 °C.
\end{itemchoice}
}
\end{ex}

% ===== Câu 50 =====
% ID: L12C1B1-71-TL
% Mức: NB
\begin{bt}
% ID: L12C1B1-71-TL
% Mức: NB
Giải thích vì sao thức ăn chín nhanh trong nồi áp suất và nêu lưu ý an toàn.
\loigiai{
Áp suất cao làm nhiệt độ sôi của nước cao hơn 100 ° $C$, nên thức ăn được nấu ở nhiệt độ cao và chín nhanh. Phải kiểm tra van, không mở nắp khi còn áp suất cao và xả áp đúng hướng dẫn.
}
\end{bt}

% ===== Câu 53 =====
% ID: L12C1B1-74-DS
% Mức: VD
\begin{ex}
% ID: L12C1B1-74-DS
% Mức: VD
Các biện pháp làm khô.
\choiceTF
{\True Phơi rộng quần áo làm tăng diện tích bay hơi.}
{\True Thông gió mang hơi ẩm đi.}
{Không khí ẩm bão hòa làm bay hơi nhanh hơn không khí khô.}
{\True Sấy ấm thường tăng tốc độ bay hơi.}
\loigiai{
\begin{itemchoice}
\item (Đúng) Phơi rộng quần áo làm tăng diện tích bay hơi.
\item (Đúng) Thông gió mang hơi ẩm đi.
\item (Sai) Không khí ẩm bão hòa làm bay hơi nhanh hơn không khí khô.
\item (Đúng) Sấy ấm thường tăng tốc độ bay hơi.
\end{itemchoice}
}
\end{ex}

% ===== Câu 54 =====
% ID: L12C1B1-78-TL
% Mức: TH
\begin{bt}
% ID: L12C1B1-78-TL
% Mức: TH
Giải thích vì sao cồn sát khuẩn trên da gây cảm giác mát và vì sao quạt làm cảm giác mát rõ hơn.
\loigiai{
Cồn bay hơi cần năng lượng nên lấy nhiệt từ da, làm da giảm nhiệt độ. Quạt mang lớp không khí giàu hơi cồn đi, tăng tốc độ bay hơi nên tốc độ lấy nhiệt tăng và cảm giác mát rõ hơn.
}
\end{bt}

% ===== Câu 57 =====
% ID: L12C1B1-79-DS
% Mức: VDC
\begin{ex}
% ID: L12C1B1-79-DS
% Mức: VDC
Về bỏng hơi nước.
\choiceTF
{\True Hơi 100 ° $C$ có thể ngưng tụ trên da.}
{\True Ngưng tụ tỏa nhiệt Lm.}
{\True Cùng khối lượng, hơi 100 ° $C$ thường truyền nhiều nhiệt hơn nước 100 ° $C$ khi cùng hạ đến nhiệt độ da.}
{Hơi nước không gây bỏng vì là chất khí.}
\loigiai{
\begin{itemchoice}
\item (Đúng) Hơi 100 ° $C$ có thể ngưng tụ trên da.
\item (Đúng) Ngưng tụ tỏa nhiệt Lm.
\item (Đúng) Cùng khối lượng, hơi 100 ° $C$ thường truyền nhiều nhiệt hơn nước 100 ° $C$ khi cùng hạ đến nhiệt độ da.
\item (Sai) Hơi nước không gây bỏng vì là chất khí.
\end{itemchoice}
}
\end{ex}

% ===== Câu 58 =====
% ID: L12C1B1-80-TL
% Mức: TH
\begin{bt}
% ID: L12C1B1-80-TL
% Mức: TH
Phân tích ảnh hưởng của nhiệt độ, diện tích mặt thoáng, gió và độ ẩm đến tốc độ bay hơi.
\loigiai{
Nhiệt độ cao làm phân tử dễ thoát khỏi mặt thoáng; diện tích lớn tăng số phân tử có thể thoát; gió mang hơi đi; độ ẩm thấp giúp không khí nhận thêm hơi. Vì vậy tăng ba yếu tố đầu và giảm độ ẩm thường làm bay hơi nhanh.
}
\end{bt}

% ===== Câu 62 =====
% ID: L12C1B1-81-TL
% Mức: VD
\begin{bt}
% ID: L12C1B1-81-TL
% Mức: VD
Giải thích vì sao bỏng hơi nước 100 ° $C$ nguy hiểm hơn bỏng nước 100 ° $C$ với cùng khối lượng.
\loigiai{
Hơi nước khi gặp da ngưng tụ và tỏa thêm nhiệt lượng Lm, sau đó nước ngưng còn nguội từ 100 ° $C$ xuống nhiệt độ da. Nước 100 ° $C$ chỉ có phần nguội đi, nên hơi truyền năng lượng lớn hơn.
}
\end{bt}

% ===== Câu 66 =====
% ID: L12C1B1-92-TL
% Mức: VD
\begin{bt}
% ID: L12C1B1-92-TL
% Mức: VD
Một người nói “phơi trong phòng kín thật nóng luôn nhanh khô nhất”. Hãy nhận xét.
\loigiai{
Không hoàn toàn đúng. Nhiệt độ cao thúc đẩy bay hơi, nhưng phòng kín nhanh trở nên ẩm, làm tốc độ bay hơi giảm. Cần thêm thông gió hoặc hút ẩm để mang hơi nước đi.
}
\end{bt}

% ===== Câu 70 =====
% ID: L12C1B1-93-TL
% Mức: VDC
\begin{bt}
% ID: L12C1B1-93-TL
% Mức: VDC
Trình bày vai trò của hóa hơi và ngưng tụ trong chu trình làm lạnh của điều hòa.
\loigiai{
Ở dàn lạnh, môi chất áp suất thấp bay hơi và hấp thụ nhiệt trong phòng. Máy nén thực hiện công; ở dàn nóng, môi chất ngưng tụ và tỏa nhiệt ra ngoài. Van tiết lưu hạ áp để chu trình lặp lại.
}
\end{bt}

% ===== Câu 88 =====
% ID: L12C1B1-94-TN
% Mức: NB
\begin{ex}
% ID: L12C1B1-94-TN
% Mức: NB
Cồn bôi lên da tạo cảm giác mát chủ yếu vì
\choice
{cồn đông đặc}
{\True cồn bay hơi thu nhiệt từ da}
{cồn ngưng tụ tỏa nhiệt}
{nhiệt độ da tự tăng}
\loigiai{
Bay hơi là quá trình thu nhiệt; cồn lấy nhiệt từ da nên da mát.
}
\end{ex}

% ===== Câu 92 =====
% ID: L12C1B1-95-TN
% Mức: NB
\begin{ex}
% ID: L12C1B1-95-TN
% Mức: NB
Trong tủ lạnh, môi chất lạnh làm lạnh khoang khi
\choice
{\True bay hơi và thu nhiệt}
{ngưng tụ và thu nhiệt}
{đông đặc và thu nhiệt}
{tăng áp không trao đổi nhiệt}
\loigiai{
Môi chất bay hơi ở dàn lạnh, hấp thụ nhiệt từ khoang.
}
\end{ex}

% ===== Câu 96 =====
% ID: L12C1B1-96-TN
% Mức: NB
\begin{ex}
% ID: L12C1B1-96-TN
% Mức: NB
Hơi nước ngưng tụ trên thành cốc lạnh làm thành cốc
\choice
{\True nhận nhiệt từ hơi}
{không trao đổi nhiệt}
{chỉ giảm khối lượng}
{hóa hơi nhanh hơn}
\loigiai{
Khi ngưng tụ, hơi nước tỏa nhiệt cho thành cốc và môi trường.
}
\end{ex}

% ===== Câu 100 =====
% ID: L12C1B1-98-TN
% Mức: TH
\begin{ex}
% ID: L12C1B1-98-TN
% Mức: TH
Quần áo nhanh khô hơn khi có gió vì gió
\choice
{làm giảm diện tích mặt thoáng}
{\True mang lớp hơi ẩm sát bề mặt đi}
{làm nước ngưng tụ}
{làm $L$ của nước bằng 0}
\loigiai{
Gió làm giảm độ ẩm cục bộ gần bề mặt, tăng tốc độ bay hơi.
}
\end{ex}

% ===== Câu 104 =====
% ID: L12C1B1-99-TN
% Mức: TH
\begin{ex}
% ID: L12C1B1-99-TN
% Mức: TH
Trải nước lên sân nóng giúp sân mát đi vì nước
\choice
{\True bay hơi thu nhiệt}
{ngưng tụ tỏa nhiệt}
{làm tăng bức xạ Mặt Trời}
{có nhiệt độ sôi bằng 0 ° $C$}
\loigiai{
Nước bay hơi lấy nhiệt từ sân và không khí gần sân.
}
\end{ex}

% ===== Câu 108 =====
% ID: L12C1B1-101-TN
% Mức: TH
\begin{ex}
% ID: L12C1B1-101-TN
% Mức: TH
Trong điều hòa không khí, bộ phận đặt trong phòng cần tạo điều kiện để môi chất
\choice
{\True bay hơi}
{ngưng tụ}
{đông đặc}
{nóng chảy}
\loigiai{
Tại dàn lạnh, môi chất bay hơi và hấp thụ nhiệt trong phòng.
}
\end{ex}

% ===== Câu 112 =====
% ID: L12C1B1-102-TN
% Mức: VD
\begin{ex}
% ID: L12C1B1-102-TN
% Mức: VD
Nồi áp suất làm thức ăn chín nhanh vì
\choice
{áp suất giảm, nhiệt độ sôi giảm}
{\True áp suất tăng, nhiệt độ sôi tăng}
{nước không cần thu nhiệt}
{nhiệt hóa hơi bằng 0}
\loigiai{
Áp suất cao làm nhiệt độ sôi của nước tăng, thức ăn được nấu ở nhiệt độ cao hơn.
}
\end{ex}

% ===== Câu 116 =====
% ID: L12C1B1-103-TN
% Mức: VD
\begin{ex}
% ID: L12C1B1-103-TN
% Mức: VD
Muốn tăng tốc độ bay hơi của nước, biện pháp phù hợp là
\choice
{giảm diện tích mặt thoáng}
{tăng độ ẩm không khí}
{\True tăng nhiệt độ và thông gió}
{đậy kín bình}
\loigiai{
Nhiệt độ cao, diện tích thoáng lớn và gió thúc đẩy bay hơi.
}
\end{ex}

% ===== Câu 120 =====
% ID: L12C1B1-104-TN
% Mức: VDC
\begin{ex}
% ID: L12C1B1-104-TN
% Mức: VDC
Bỏng do hơi nước 100 ° $C$ thường nặng hơn do nước 100 ° $C$ vì hơi nước
\choice
{có khối lượng bằng 0}
{\True khi ngưng tụ còn tỏa nhiệt Lm}
{không truyền nhiệt}
{làm da bay hơi tức thì}
\loigiai{
Hơi ngưng tụ trên da tỏa thêm nhiệt hóa hơi lớn.
}
\end{ex}

% ===== Câu 123 =====
% ID: L12C1B1-105-TN
% Mức: VDC
\begin{ex}
% ID: L12C1B1-105-TN
% Mức: VDC
Sự bay hơi khác sự sôi ở chỗ bay hơi
\choice
{chỉ xảy ra ở nhiệt độ sôi}
{\True xảy ra trên mặt thoáng ở mọi nhiệt độ}
{xảy ra trong toàn khối ở nhiệt độ xác định}
{không thu nhiệt}
\loigiai{
Bay hơi xảy ra trên mặt thoáng và có thể ở mọi nhiệt độ.
}
\end{ex}
